\section{Introduction}
\label{sec:intro}

A backdoored classifier behaves normally on clean inputs and predicts an
attacker's class whenever a trigger appears in the input
\cite{li2022backdoorsurvey}. The attacker plants the behavior by poisoning a
small share of the training data, often 1\% or less. The defender receives the
trained model without ever seeing that data. Input-level detection asks the
most practical question in this setting. Given the deployed model and 1 incoming
input, is that input carrying a trigger? A detector that answers well needs no
retraining and runs in front of the deployed model.

Prediction Shift Backdoor Detection (PSBD) \cite{li2025psbd} is such a detector.
Its observation is that a backdoored input is decided by a shortcut the model
learned quickly and firmly, while a clean input is decided by a broad set of
features. When the network is perturbed at inference the broad decision wobbles
and the shortcut holds. PSBD therefore runs the model a few times with dropout
active, measures how much the probability of the predicted class falls and flags
the inputs whose prediction barely moves. It needs only a small clean validation
set to set its threshold, and on ResNet-18 and VGG-16 it reports strong results
across 7 attacks.

Nothing in that argument depends on convolutions, yet the perturbation does not
carry over unchanged. A ResNet basic block has 1 natural place for dropout, the
residual stream after the add, and the original work put it there. A transformer
block has no such default. The input to attention, the input to the MLP, each
branch output before its residual add, the stream after each add, the attention
heads and the MLP hidden units are all different tensors with different roles.
Nor is dropout the only thing to inject, since whole tokens or channels can be
masked and noise can be added.

\paragraph{Research questions.} Adapting PSBD to a Vision Transformer (ViT)
\cite{dosovitskiy2021an} is therefore a design problem on 2 axes, where to
perturb and what to inject. We ask whether the choice matters, whether the 2
axes act separately, and why the winning placement wins, since a ranking no
mechanism explains predicts nothing about the next architecture.

The choice decides the outcome. Masking whole tokens at the input of every
attention block, which we call PSBD-TM, reaches a mean AUROC of
\HeadlineAurocAdaptive{} against \PublishedAurocAdaptive{} for our reading of the
original placement, which we call PSBD-RD. Site and operator both move the score,
and what decides the operator is the side of the LayerNorm the perturbation lands
on rather than the site. The mechanism is that a backdoor in a ViT is 1 direction
in the residual stream, written in the last third of the network and routed
through attention from the trigger's patch tokens to the class token.

\paragraph{Contributions.} We make 3 claims of different kinds.

\begin{enumerate}
  \item A systematic result. Over \PanelCellsCached{} backdoored ViT-B/16 models on
    \PanelDatasetsCached{} datasets, PSBD-TM beats PSBD-RD by
    \HeadlineGainAdaptiveAuroc{} [\HeadlineGainAdaptiveAurocLow{},
    \HeadlineGainAdaptiveAurocHigh{}], with its largest advantage at \PanelRateLow{}
    poisoning, and by \SwinGainRecommendedMinusPublished{} over \SwinCells{} Swin-S
    models (\cref{sec:results,sec:swin}). Against \DetectorsComparedCount{} ported
    input-level detectors on the same splits it ranks \DetectorsAurocRankOurs{} of
    \DetectorsAurocDefensesRanked{} defenses where PSBD-RD ranks
    \DetectorsAurocRankPublished{}, though its margin over the strongest competitor is
    only \DetectorsAurocMargin{} (\cref{sec:robustness:detectors}).
  \item 2 refutations. The claim that dropout before the residual add beats
    dropout after it, which motivated this project, does not hold:
    \PreMinusPostMatched{} [\PreMinusPostMatchedLow{},
    \PreMinusPostMatchedHigh{}] at matched disturbance over \PreMinusPostN{}
    models. Nor is the winner the input-side family of sites, since the family
    gap changes sign with the operator (\cref{sec:results:family}).
  \item A mechanistic account, tested on the same models. It explains the winning
    site, the choice of whole tokens over coordinates and the reversal through
    LayerNorm absorption (\cref{sec:mechanism}).
\end{enumerate}

Every measurement in this paper is produced by a script from a coverage ledger
that declares the models before any sweep, and the build fails on a hand-typed one.
The code, the ledger and the generators are released. \Cref{sec:limitations} states
what the evidence does not cover.
